// 323111047

\section{Introducción}

Planteamiento del problema. El reto técnico de esta práctica consiste en desarrollar una aplicación de consola en lenguaje C que simule un carrusel publicitario digital mediante el uso de una lista doble circular ligada. El problema central radica en la gestión de nodos que posean enlaces simultáneos al elemento anterior y al siguiente, asegurando que el primer nodo esté conectado con el último en ambos sentidos de manera permanente. Se deben implementar funciones para registrar, buscar y eliminar anuncios por identificador, además de permitir una navegación bidireccional infinita que mantenga la integridad de la memoria dinámica al modificar la estructura de la lista durante la ejecución.

\vspace{0.4cm}

Motivación. La necesidad de dar solución a este problema se basa en la importancia de dominar estructuras de datos que permitan una navegación flexible y eficiente en aplicaciones de software reales. En el desarrollo de interfaces de usuario modernas, los carruseles de contenido son herramientas fundamentales que requieren un ciclo permanente de datos para mejorar la interacción con el usuario. Estudiar las listas doblemente circulares permite al estudiante comprender cómo optimizar el acceso a la información en ambos sentidos y cómo administrar punteros complejos, lo cual es esencial para el diseño de sistemas de archivos y gestores de memoria avanzados en el ámbito profesional.

\vspace{0.4cm}

Objetivos. Se busca que el alumno demuestre dominio en la manipulación de enlaces anteriores y siguientes dentro de un nodo, validando que las operaciones de inserción y borrado mantengan siempre la circularidad bidireccional de la estructura. Finalmente, se pretende verificar que el programa permita una navegación fluida hacia el anuncio siguiente y el anterior de manera infinita.